\documentclass[11pt]{article}
\usepackage[a4paper,margin=2.2cm]{geometry}
\usepackage[T1]{fontenc}
\usepackage{lmodern}
\usepackage{url}
\setlength{\parindent}{0pt}
\setlength{\parskip}{6pt}
\pagestyle{empty}

\begin{document}

\begin{flushright}
Swati Hira\\
Department of Computer Science and Engineering\\
IIIT Nagpur, Nagpur, India\\
\url{shira@iiitn.ac.in}\\[4pt]
\today
\end{flushright}

The Editor-in-Chief\\
\textit{Remote Sensing Letters}\\
Taylor \& Francis

Dear Editor,

We submit the manuscript \textbf{``Channel attention in lightweight hyperspectral CNNs: an equivalence-tested comparison''} for consideration as a Research Letter in \textit{Remote Sensing Letters}.

Channel-attention blocks are frequently proposed for lightweight hyperspectral classifiers, and their gains are usually reported on random pixel splits in which overlapping patches leak spatial context into the test set. Leakage itself has been documented for whole architectures; whether the \emph{choice of block} still matters once evaluation is tightened has, to our knowledge, not been tested with paired statistics. Our letter does this:

\begin{itemize}\setlength{\itemsep}{2pt}
\item Seven channel-modulation blocks (SE, ECA, CBAM, GCT, SRM and two variants of a dispersion-aware probe, SNSM) are compared on one backbone with identical inputs and schedule, no test-based model selection and up to fifteen paired seeds, on five scenes under seven protocols that vary class counts and spatial overlap, with checks of host depth, learning rate and feature fitting. The absence of a gain is tested directly with equivalence tests ($\pm$1 point of average accuracy), not inferred from non-significance.
\item Once every class has at least 15 training pixels, all blocks are equivalent within $\pm$1 point. The one large difference, 2--4 points for dispersion-based gates on random Indian Pines splits, comes from two single-pixel classes, is surpassed by a deeper plain CNN, shrinks at a higher learning rate, does not reappear when the same scarcity is imposed on Pavia University or Kennedy Space Center, and is mostly not significant on patch-disjoint splits, where every block loses more than 30 points.
\item From this we derive a simple check for future block proposals: report the training pixels of the smallest classes, test for equivalence, and repeat the comparison on patch-disjoint splits.
\end{itemize}

We believe these results suit the journal's readership because they bear directly on how new modules for hyperspectral classification are evaluated and compared. The findings are short, self-contained and supported by complete code and per-run results, which are publicly available at \url{https://github.com/Devguru-codes/research_paper_code}.

The manuscript is original, has not been published previously and is not under consideration elsewhere. All authors have approved the submission. The authors declare no competing interests, and the research received no external funding. Generative AI was used to assist with experiment code, reference checking and language editing, as stated in the manuscript; the authors verified all content and take full responsibility for it.

Thank you for considering our work.

Yours sincerely,

Swati Hira (corresponding author), on behalf of Devguru Tiwari, Pratibha Kokardekar and S.~R. Sathe

\end{document}
